
\documentclass[11pt,a4paper]{article}
\usepackage[utf8]{inputenc}
\usepackage[T1]{fontenc}
\usepackage[french]{babel}
\usepackage{lmodern}
\usepackage{microtype}
\usepackage[a4paper,margin=1.7cm,headheight=15pt]{geometry}
\usepackage{amsmath,amssymb,mathtools}
\usepackage{array,tabularx,multicol}
\usepackage{enumitem}
\usepackage[most]{tcolorbox}
\usepackage{xcolor}
\usepackage{tikz}
\usepackage{pgfplots}
\pgfplotsset{compat=1.18}
\usetikzlibrary{arrows.meta,calc,patterns}
\usepackage{fancyhdr}
\usepackage{fancyvrb}
\usepackage{listings}
\usepackage{hyperref}
\usepackage{pdfpages}

\definecolor{bleufonce}{RGB}{28,55,120}
\definecolor{bleuclair}{RGB}{238,243,255}
\definecolor{grisclair}{RGB}{246,246,246}
\definecolor{tableline}{RGB}{45,45,45}
\definecolor{softgray}{RGB}{246,247,249}
\definecolor{deepblue}{RGB}{25,62,115}
\definecolor{softblue}{RGB}{235,241,250}
\definecolor{accent}{RGB}{20,82,130}

\hypersetup{colorlinks=true,linkcolor=bleufonce,urlcolor=bleufonce,pdfauthor={},pdftitle={Parcours de révision -- Trigonométrie et LGN -- Corrigés}}
\setlength{\parindent}{0pt}
\setlength{\parskip}{0.28em}
\renewcommand{\arraystretch}{1.12}
\setlist[itemize]{leftmargin=1.25em,itemsep=0.15em,topsep=0.2em}
\setlist[enumerate,1]{label=\textbf{\arabic*.}, leftmargin=1.05cm, itemsep=0.35em, topsep=0.25em}
\setlist[enumerate,2]{label=\textbf{\alph*.}, leftmargin=0.85cm, itemsep=0.25em, topsep=0.2em}

\pagestyle{fancy}
\fancyhf{}
\lhead{}
\rhead{}
\cfoot{\thepage}

\newcommand{\e}{\mathrm e}
\newcommand{\dd}{\,\mathrm d}
\newcommand{\vect}[1]{\overrightarrow{#1}}
\newcommand{\origine}[1]{ {\scriptsize\textcolor{bleufonce}{(site #1)}}}
\newcommand{\code}[1]{\texttt{\detokenize{#1}}}
\newcommand{\Bin}{\mathcal B}
\newcounter{question}
\newcounter{reponse}

\newenvironment{blocquestions}[1]{%
  \begin{tcolorbox}[
    enhanced,
    breakable,
    colback=bleuclair,
    colframe=bleufonce,
    coltitle=white,
    colbacktitle=bleufonce,
    title={\bfseries #1},
    fonttitle=\bfseries,
    boxrule=0.6pt,
    arc=2mm,
    left=2mm,
    right=2mm,
    top=2mm,
    bottom=1.5mm,
    before skip=3mm,
    after skip=3mm, fontupper=\large
  ]
  \large
  \begin{enumerate}[leftmargin=*,label=\textbf{\arabic*.},itemsep=0.45em,topsep=0.4em]
  \setcounter{enumi}{\value{question}}
}{%
  \setcounter{question}{\value{enumi}}
  \end{enumerate}
  \end{tcolorbox}
}

\newenvironment{blocreponses}[1]{%
  \begin{tcolorbox}[
    enhanced,
    breakable,
    colback=grisclair,
    colframe=bleufonce,
    coltitle=white,
    colbacktitle=bleufonce,
    title={\bfseries #1},
    fonttitle=\bfseries,
    boxrule=0.6pt,
    arc=2mm,
    left=2mm,
    right=2mm,
    top=2mm,
    bottom=1.5mm,
    before skip=3mm,
    after skip=3mm, fontupper=\large
  ]
  \large
  \begin{enumerate}[leftmargin=*,label=\textbf{\arabic*.},itemsep=0.45em,topsep=0.5em]
  \setcounter{enumi}{\value{reponse}}
}{%
  \setcounter{reponse}{\value{enumi}}
  \end{enumerate}
  \end{tcolorbox}
}

\newtcolorbox{correctionbox}[1]{
  enhanced,
  breakable,
  colback=grisclair,
  colframe=black!55,
  boxrule=0.55pt,
  arc=2mm,
  left=2mm,right=2mm,top=1.5mm,bottom=1.5mm,
  title=\textbf{#1},
  coltitle=black,
  colbacktitle=black!12,
  before skip=2mm,
  after skip=2mm
}

\newcommand{\titreSujet}[2]{%
  \subsection*{#1}
  \addcontentsline{toc}{subsection}{#1}
  \textit{Référence : #2}\par\medskip\hrule\medskip
}

\lstset{language=Python,basicstyle=\ttfamily\small,numbers=left,numberstyle=\tiny,frame=single,showstringspaces=false,columns=fullflexible,keepspaces=true}
\sloppy

\begin{document}
\begin{center}
{\LARGE\bfseries Parcours de révision -- Trigonométrie et loi des grands nombres}\\[1mm]

\end{center}
\vspace{1mm}
\tableofcontents
\clearpage

\section*{Partie 1 -- Corrections des questions flashs}
\addcontentsline{toc}{section}{Partie 1 -- Corrections des questions flashs}
\setcounter{reponse}{0}
\begin{blocreponses}{Réponses -- A. Trigonométrie : cercle et propriétés}
  \item \origine{479} C'est le cercle de centre $O$ et de rayon $1$, orienté dans le sens direct.
  \item \origine{480} Cela signifie que, pour tout réel $x$, $\cos(x+2\pi)=\cos(x)$ et $\sin(x+2\pi)=\sin(x)$.
  \item \origine{481} La fonction cosinus est paire : pour tout réel $x$, $\cos(-x)=\cos(x)$.
  \item \origine{482} La fonction sinus est impaire : pour tout réel $x$, $\sin(-x)=-\sin(x)$.
  \item \origine{528} C'est le point de coordonnées $(\cos x\,;\,\sin x)$.
  \item \origine{529} Un tour complet correspond à une longueur $2\pi$.
  \item \origine{530} On a toujours $\cos^2 x+\sin^2 x=1$.
\end{blocreponses}
\begin{blocreponses}{Réponses -- B. Trigonométrie : valeurs remarquables}
  \item \origine{531} On a $\cos(0)=1$ et $\sin(0)=0$.
  \item \origine{533} On a $\cos\left(\dfrac{\pi}{2}\right)=0$ et $\sin\left(\dfrac{\pi}{2}\right)=1$.
  \item \origine{484} On a $\cos\left(\dfrac{\pi}{4}\right)=\dfrac{\sqrt2}{2}$ et $\sin\left(\dfrac{\pi}{4}\right)=\dfrac{\sqrt2}{2}$.
  \item \origine{691} On a $\cos\left(\dfrac{\pi}{6}\right)=\dfrac{\sqrt3}{2}$ et $\sin\left(\dfrac{\pi}{6}\right)=\dfrac12$.
\end{blocreponses}
\begin{blocreponses}{C. Liens avec l'analyse}
  \item \origine{158} Une primitive est $\sin(x)$.
  \item \origine{159} Une primitive est $-\cos(x)$.
  \item \origine{247} Une primitive est $\cos x$.
  \item \origine{798} Valeur classique : la limite vaut $\dfrac12$.
\end{blocreponses}
\begin{blocreponses}{Réponses -- D. Variables aléatoires et loi des grands nombres}
  \item Si $X$ prend les valeurs $x_1,\ldots,x_n$, alors $E(X)=\displaystyle\sum_{i=1}^n x_iP(X=x_i)$.
  \item $V(X)=E\bigl((X-E(X))^2\bigr)$, ou $V(X)=\displaystyle\sum_i (x_i-E(X))^2P(X=x_i)$.
  \item $\sigma(X)=\sqrt{V(X)}$.
  \item $E(X)=np$.
  \item $V(X)=np(1-p)$.
  \item Elle représente une valeur moyenne attendue sur un grand nombre de répétitions.
  \item Elle représente la moyenne observée sur l'échantillon.
  \item Lorsque le nombre de répétitions devient grand, la moyenne observée se rapproche de l'espérance.
  \item Il diminue.
  \item Par linéarité de l'espérance, $E(M_n)=\mu$.
\end{blocreponses}

\clearpage
\section*{Partie 2 -- Corrigés des sujets de bac}
\addcontentsline{toc}{section}{Partie 2 -- Corrigés des sujets de bac}
\titreSujet{Sujet 1 -- Loi des grands nombres et moyenne d'échantillons}{Métropole, juin 2025, sujet 1, exercice 1}

\begin{correctionbox}{1. Arbre pondéré}
On complète d'abord les probabilités des quatre groupes sanguins :
\[
p(O)=1-p(A)-p(B)-p(AB)=1-0{,}45-0{,}10-0{,}03=0{,}42.
\]
Les probabilités conditionnelles contraires valent :
\[
p_A(\overline R)=0{,}15,\qquad p_B(\overline R)=0{,}16,\qquad p_{AB}(\overline R)=0{,}18.
\]

\begin{center}
\begin{tikzpicture}[grow=right, level distance=2.8cm, sibling distance=18mm,
  every node/.style={font=\small}, edge from parent/.style={draw},
  edge from parent path={(\tikzparentnode.east) -- (\tikzchildnode.west)}]
\node {}
  child {node {$O$}
    child {node {$\overline R$}}
    child {node {$R$}}
    edge from parent node[below] {$0{,}42$}}
  child {node {$AB$}
    child {node {$\overline R$} edge from parent node[below] {$0{,}18$}}
    child {node {$R$} edge from parent node[above] {$0{,}82$}}
    edge from parent node[below] {$0{,}03$}}
  child {node {$B$}
    child {node {$\overline R$} edge from parent node[below] {$0{,}16$}}
    child {node {$R$} edge from parent node[above] {$0{,}84$}}
    edge from parent node[above] {$0{,}10$}}
  child {node {$A$}
    child {node {$\overline R$} edge from parent node[below] {$0{,}15$}}
    child {node {$R$} edge from parent node[above] {$0{,}85$}}
    edge from parent node[above] {$0{,}45$}};
\end{tikzpicture}
\end{center}
\end{correctionbox}

\begin{correctionbox}{2. Calcul de $p(B\cap R)$}
On utilise la formule d'une probabilité d'intersection avec une probabilité conditionnelle :
\[
p(B\cap R)=p(B)\times p_B(R)=0{,}10\times 0{,}84=0{,}084.
\]
Ainsi, la probabilité qu'une personne choisie au hasard soit de groupe sanguin B et de rhésus positif est égale à $0{,}084$, soit $8{,}4\,\%$.
\end{correctionbox}

\begin{correctionbox}{3. Calcul de $p_O(R)$}
Les événements $A$, $B$, $AB$ et $O$ forment une partition de l'univers. D'après la formule des probabilités totales, on a :
\[
p(R)=p(A\cap R)+p(B\cap R)+p(AB\cap R)+p(O\cap R).
\]
Donc
\[
0{,}8397=0{,}45\times0{,}85+0{,}10\times0{,}84+0{,}03\times0{,}82+0{,}42\times p_O(R).
\]
Or
\[
0{,}45\times0{,}85+0{,}10\times0{,}84+0{,}03\times0{,}82=0{,}4911.
\]
Ainsi,
\[
0{,}42\times p_O(R)=0{,}8397-0{,}4911=0{,}3486,
\]
puis
\[
p_O(R)=\dfrac{0{,}3486}{0{,}42}=0{,}83.
\]
\end{correctionbox}

\begin{correctionbox}{4. Probabilité d'être donneur universel}
Un donneur universel est une personne de groupe O et de rhésus négatif. On cherche donc $p(O\cap \overline R)$.

D'après la question précédente, $p_O(R)=0{,}83$, donc $p_O(\overline R)=0{,}17$.

Ainsi,
\[
p(O\cap \overline R)=p(O)\times p_O(\overline R)=0{,}42\times0{,}17=0{,}0714.
\]
La probabilité qu'un individu choisi au hasard dans la population française soit donneur universel est donc $0{,}0714$, soit $7{,}14\,\%$.
\end{correctionbox}

\begin{correctionbox}{5. Loi binomiale}
\begin{enumerate}
  \item On répète $100$ expériences de Bernoulli indépendantes et identiques, car le choix de l'échantillon est assimilé à un tirage avec remise. Le succès est : « la personne choisie est donneur universel », de probabilité $p=0{,}0714$. La variable aléatoire $X$ compte le nombre de succès parmi les 100 personnes. Donc
  \[
  X\sim \mathcal B(100;0{,}0714).
  \]

  \item On cherche $P(X\leqslant 7)$. À la calculatrice, avec $X\sim\mathcal B(100;0{,}0714)$,
  \[
  P(X\leqslant 7)\approx 0{,}577.
  \]
  La probabilité qu'il y ait au plus 7 donneurs universels dans cet échantillon est donc environ $0{,}577$.

  \item Pour une loi binomiale $\mathcal B(n;p)$, on a $E(X)=np$ et $V(X)=np(1-p)$. Donc
  \[
  E(X)=100\times0{,}0714=7{,}14,
  \]
  et
  \[
  V(X)=100\times0{,}0714\times(1-0{,}0714)=6{,}630204\approx 6{,}63.
  \]
\end{enumerate}
\end{correctionbox}

\begin{correctionbox}{6. Moyenne sur $N$ villes et Bienaymé-Tchebychev}
\begin{enumerate}
  \item La variable aléatoire $M_N$ représente le nombre moyen de donneurs universels dans les échantillons de 100 personnes prélevés dans les $N$ villes.

  \item Par linéarité de l'espérance,
  \[
  E(M_N)=E\left(\dfrac{X_1+X_2+\cdots+X_N}{N}\right)
  =\dfrac{E(X_1)+E(X_2)+\cdots+E(X_N)}{N}.
  \]
  Comme chaque variable $X_i$ a pour espérance $7{,}14$, on obtient
  \[
  E(M_N)=\dfrac{N\times7{,}14}{N}=7{,}14.
  \]

  \item Comme les variables $X_1,\ldots,X_N$ sont indépendantes,
  \[
  V(X_1+\cdots+X_N)=V(X_1)+\cdots+V(X_N)=N\times6{,}63.
  \]
  Donc
  \[
  V(M_N)=V\left(\dfrac{X_1+\cdots+X_N}{N}\right)=\dfrac{1}{N^2}V(X_1+\cdots+X_N)
  =\dfrac{N\times6{,}63}{N^2}=\dfrac{6{,}63}{N}.
  \]

  \item On remarque que
  \[
  7<M_N<7{,}28 \quad \Longleftrightarrow \quad |M_N-7{,}14|<0{,}14.
  \]
  D'après l'inégalité de Bienaymé-Tchebychev,
  \[
  P\left(|M_N-E(M_N)|\geqslant 0{,}14\right)\leqslant \dfrac{V(M_N)}{0{,}14^2}.
  \]
  Donc
  \[
  P\left(|M_N-7{,}14|\geqslant 0{,}14\right)\leqslant \dfrac{6{,}63}{N\times0{,}0196}.
  \]
  Ainsi,
  \[
  P(7<M_N<7{,}28)\geqslant 1-\dfrac{6{,}63}{N\times0{,}0196}.
  \]
  On veut que ce minorant soit strictement supérieur à $0{,}95$ :
  \[
  1-\dfrac{6{,}63}{N\times0{,}0196}>0{,}95.
  \]
  Cela équivaut à
  \[
  \dfrac{6{,}63}{N\times0{,}0196}<0{,}05,
  \]
  donc
  \[
  N>\dfrac{6{,}63}{0{,}05\times0{,}0196}\approx6765{,}31.
  \]
  La plus petite valeur entière possible est donc
  \[
  \boxed{N=6766}.
  \]
\end{enumerate}
\end{correctionbox}
\clearpage
\titreSujet{Sujet 2 -- Trigonométrie, convexité et intégration}{Amérique du Nord, 22 mai 2025, sujet 2, exercice 4}

\begin{correctionbox}{Partie A}
\begin{enumerate}
\item On dérive :
\[
f'(x)=\e^x\sin x+\e^x\cos x=\e^x(\sin x+\cos x).
\]
Sur $\left[0;\frac\pi2\right]$, $\sin x\geqslant0$ et $\cos x\geqslant0$, et leur somme est strictement positive. Donc $f'(x)>0$ et $f$ est strictement croissante sur cet intervalle.

\item On a $f(0)=0$ et $f'(0)=1$. La tangente en $0$ a donc pour équation
\[
y=x.
\]
Ensuite
\[
f''(x)=\e^x(\sin x+\cos x)+\e^x(\cos x-\sin x)=2\e^x\cos x.
\]
Sur $\left[0;\frac\pi2\right]$, $f''(x)\geqslant0$, donc $f$ est convexe.
Une fonction convexe est au-dessus de ses tangentes ; en particulier, au-dessus de la tangente en $0$. Donc
\[
\e^x\sin x\geqslant x
\]
pour tout $x\in\left[0;\frac\pi2\right]$.

\item Le signe de $f''(x)=2\e^x\cos x$ change lorsque $x$ traverse $\frac\pi2$ : positif avant, négatif après. Le point d'abscisse $\frac\pi2$ est donc un point d'inflexion.
\end{enumerate}
\end{correctionbox}

\begin{correctionbox}{Partie B}
\begin{enumerate}
\item Première intégration par parties, avec $u=\sin x$ et $v'=\e^x$ :
\[
I=\left[\e^x\sin x\right]_0^{\pi/2}-\int_0^{\pi/2}\e^x\cos x\,dx
=\e^{\pi/2}-J.
\]
Deuxième intégration par parties, avec $u=\e^x$ et $v'=\sin x$ :
\[
I=\left[-\e^x\cos x\right]_0^{\pi/2}+\int_0^{\pi/2}\e^x\cos x\,dx
=1+J.
\]

\item En additionnant les deux expressions,
\[
2I=1+\e^{\pi/2},
\]
donc
\[
I=\dfrac{1+\e^{\pi/2}}2.
\]

\item Sur $\left[0;\frac\pi2\right]$, on a $f(x)\geqslant x$. L'aire vaut donc
\[
\int_0^{\pi/2}\left(\e^x\sin x-x\right)\,dx
=I-\left[\dfrac{x^2}{2}\right]_0^{\pi/2}.
\]
Ainsi
\[
\mathcal A=\dfrac{1+\e^{\pi/2}}2-\dfrac{\pi^2}{8}.
\]
\end{enumerate}
\end{correctionbox}

\newpage
\end{document}
